\section{Matched-support identification experiments}\label{sec:identification}
\subsection{Observations, fixed cases and computational scale}
The point-map experiments use only the unchanged Brazilian rectangle, 2015 and 2023 native AGB, WGS84 cell areas, Collection 11 area fractions and integer overlap certificates described above. The footprint contains 65,536 cells and 61,830.3436 ha. All cells, including observed zeros, mixed cells and changing labels, are retained. The spatial experiment compares sides $s\in\{1,2,4,8,16,32,64,128,256\}$ native cells. It holds the complete fine-grid group geometry fixed while replacing observed cell masses by their sums within aligned blocks. Thus it isolates loss of mass-allocation information rather than a change in case selection or mapped-group definition. It is not a model of the sensor's point-spread function, an effective-resolution estimate or a new source observation. The analyses are exploratory; the footprint and endpoints were fixed by the earlier geometry-only support rule, not selected using these results.

The full-cell masses are $463{,}922.52584046$ and $529{,}164.8923703252$ Mg. The native-cell positive and negative point-map differences total $142{,}213.3577547465$ and $76{,}970.9912248813$ Mg. Their net is $65{,}242.3665298653$ Mg and their cancellation fraction is $0.7023402135$. These values are not adjusted for observation error. The same net is preserved at all nine scales, up to floating-point summation. At side 64, 128 and 256, the blockwise negative sum and cancellation fraction are zero. No conclusion about true fine-scale change follows without a compatible error model.

\subsection{Sharp block bounds and nested information sets}
Within block $b$, the observed mass is $q_b$ and the mapped group either has zero area, occupies the whole block, or has positive area together with a positive-area complement. Under nonnegative unconstrained density, the group mass is respectively $0$, $q_b$, or any value in $[0,q_b]$. The mixed case is attained by assigning component means $x/A_g$ and $(q_b-x)/A_h$ for arbitrary $x\in[0,q_b]$, where $A_g,A_h>0$. Summing these independent intervals yields the block endpoint bounds in the main text; taking the cross-endpoint extrema yields sharp difference bounds. The union's topology is certified absent only if every constituent count is zero, and pure only if every group count equals its total count. Counts, not rounded area fractions, determine these categories.

Every assignment satisfying finer block totals satisfies the coarser totals obtained by summation. Hence the projected group-mass sets can only expand under nested coarsening. In this experiment the group and complement each occupy positive area at both endpoints, giving the one-window difference set $[-Q_{2015},Q_{2023}]$. Its width is $993{,}087.4182107851$ Mg. The native-cell widths are $505{,}973.6696403596$ and $437{,}643.5603800343$ Mg, reductions of 49.0504\% and 55.9310\%, respectively (rounded values in the main table). Width is an identification metric conditional on fixed point maps; its reduction is not empirical accuracy or confidence-coverage improvement.

For cancellation, the triangle inequality gives $\sum_b|\sum_{i\in b}\Delta q_i|\le\sum_i|\Delta q_i|$, while the net is unchanged. At a nonzero net, $K=1-|\sum_i\Delta q_i|/\sum_i|\Delta q_i|$ therefore cannot increase under nested aggregation. If the total variation is zero, $K$ is undefined rather than assigned zero. The present nonzero-net data do not encounter that undefined case.

\subsection{Constructive endpoint witnesses}
Let $c_{it}^g$ and $n_i$ be group and total intersection counts. Define $x_{it}^{\rm lo}=q_{it}\mathbf1\{c_{it}^g=n_i\}$ and $x_{it}^{\rm hi}=q_{it}\mathbf1\{c_{it}^g>0\}$. The minimum-change witness uses $(x_{i,2015}^{\rm hi},x_{i,2023}^{\rm lo})$ and the maximum uses $(x_{i,2015}^{\rm lo},x_{i,2023}^{\rm hi})$. For each endpoint and cell, group and complement masses are nonnegative, sum to the unchanged observed cell mass, and respect absent/pure topology. The component densities in mixed cells are formed on their positive mapped areas, so the witness retains the exact mapped fraction rather than changing the map. No finite density cap, temporal growth restriction or ecological plausibility claim is added. Convex interpolation of the endpoint allocations attains all intermediate differences. The witnesses are constructed separately for each overlapping mapped group; simultaneous attainment of both groups' extreme bounds is not asserted.

The four aggregate witness certificates report 65,536 constraint-satisfying cells at each endpoint. The minimum and maximum changes reproduce the unrounded primary difference bounds to numerical summation precision. Native-cell masses and allocations are generated in memory for these checks but are not redistributed. Aggregate certificates are supplied as \path{data/extremal_witness_certificates_F11.csv}; the complete nested-scale results are in \path{data/resolution_frontier_F11.csv} and \path{data/cancellation_by_scale_F11.csv}.

\subsection{Sharp joint set for nested cover groups}
Let $A$ denote formations, $S$ added savanna and $O$ the complement, so $H=A+S$ denotes formations plus savanna. Their integer overlap counts are $c^A$, $c^S=c^H-c^A$ and $c^O=n-c^H$. The nesting $0\le c^A\le c^H\le n$ is checked for every cell. Given cell mass $q$, its feasible $(Q^A,Q^H)$ allocation is the convex hull of whichever of $(q,q)$, $(0,q)$ and $(0,0)$ correspond to components with positive mapped area. Cells with one available component give a point, two give an edge, and three give a triangle. Component area fractions are unchanged; the nonnegative component densities are free.

The endpoint polygon is the Minkowski sum of these cell polygons, and the difference polygon is the Minkowski sum of the 2023 endpoint polygon and the negative 2015 polygon. Every cell edge is horizontal, vertical or parallel to $(1,1)$. The support function of a Minkowski sum is the sum of support functions, so its possible facet orientations are contained in the same three directions, including their opposite normals. Its exact half-space description is therefore given by the lower and upper support bounds for $a$, $h$ and $h-a$. Independent cellwise endpoint optimization gives precisely the three univariate difference intervals in Eq.~\ref{eq:jointSI}:
\begin{equation}
\mathcal D_{\rm joint}=\{(a,h):a\in\widehat{\mathcal D}^A,\ h\in\widehat{\mathcal D}^H,\ h-a\in\widehat{\mathcal D}^S\}.\label{eq:jointSI}
\end{equation}
This establishes a sharp joint set, not merely necessary marginal constraints. It also explains why simultaneously choosing arbitrary marginal extrema need not be feasible. No cross-group statistical independence is assumed.

For computation, cells are aggregated into the seven nonempty availability patterns of the three components, retaining each pattern's total mass at each endpoint. Let $M_{pt}$ be that mass and $z_{pkt}$ the mass assigned to an available component $k$. The constraints are $z_{pkt}\ge0$, $\sum_kz_{pkt}=M_{pt}$ and $z_{pkt}=0$ for an absent component. This reduction is exact: for $M_{pt}>0$, assign each constituent cell's mass in the proportions $z_{pkt}/M_{pt}$; a zero-mass pattern requires only zero allocation. All cell totals, absent components and pure cells are then preserved. Thus the optimization does not discard native-cell mass constraints.

Four linear programs maximize a common signed margin $\epsilon\ge0$ under $\sigma_A a\ge\epsilon$ and $\sigma_H h\ge\epsilon$, for the four pairs $(\sigma_A,\sigma_H)\in\{-1,1\}^2$. All attain strictly positive margins. The coherent witness differences $(a,h)$, in $10^3$ Mg, are $(-223.97,-178.11)$, $(-119.38,119.38)$, $(103.68,-103.68)$ and $(282.00,259.53)$. Reconstructed cell allocations preserve mass to at most $1.43\times10^{-14}$ Mg per cell. This tolerance concerns numerical summation, not ecological uncertainty. The calculated joint polygon has six vertices. A separate support-function optimization over the pattern constraints agrees with the half-space polygon over 73 evaluated directions to $1.17\times10^{-10}$ Mg; the sharpness result rests on the geometric argument above, not that finite check.

Pattern summaries, quadrant certificates and polygon vertices are supplied as \path{data/joint_availability_patterns_F11.csv}, \path{data/joint_quadrant_certificates_F11.csv} and \path{data/joint_polygon_vertices_F11.csv}. They disclose aggregate derived quantities, not native raster arrays. The common ratio sensitivity in the next subsection is evaluated separately for each mapped group and is not asserted to define one jointly calibrated three-component ecological model.

\subsection{Sharp bounds under an assumed density contrast}
In a mixed cell write $d=f d_g+(1-f)d_h$, where $0<f<1$ and $q=Ad$. The restrictions $d_g\le R d_h$ and $d_h\le R d_g$ with $R\ge1$ imply
\[
\frac{f}{f+R(1-f)}\le\frac{s_g}{q}\le\frac{Rf}{Rf+1-f},\qquad q>0.
\]
The lower endpoint is attained with $d_h=R d_g$, and the upper with $d_g=R d_h$. Both endpoints respect the observed mixture. If $q=0$, nonnegativity gives $s_g=0$ without division. Pure/absent cells are fixed separately and impose no comparison with an unobserved zero-area component. These formulas use area fractions, not intersection-count fractions. Integer topology supplies exact zero/one values for certified absent/pure cells; floating area-integration deviations of at most $3.34\times10^{-16}$ at these certified endpoints are set to their exact values. Every mixed-cell fraction is retained strictly inside $(0,1)$, without a tolerance-based membership decision.

Independent summation and cross-endpoint subtraction yield sharp aggregate difference bounds for each $R$. They expand monotonically with $R$, reduce to homogeneous allocation at $R=1$, and converge to the primary unrestricted bounds as $R\to\infty$. The common bound applies to both endpoints and every mixed cell; it is an additional structural assumption, not a density parameter fitted from the maps. The strictly positive homogeneous differences lose strict sign identification when the lower bound reaches zero, at $R=1.4555078889$ for formations and $R=1.9495126065$ for formations plus savanna. A bracketed root calculation in log-$R$ space locates these boundaries. They are conditional sensitivity boundaries, not estimated ecological contrasts, confidence thresholds or universal necessary conditions under every possible auxiliary-information model.

\begin{longtable}{lrrr}
\caption{Conditional density-contrast sensitivity, in $10^3$ Mg. $R$ is assumed, not estimated or calibrated. $R=1$ is homogeneous allocation; $R=\infty$ is the primary unrestricted model.}\label{tab:contrast}\\
\toprule
Mapped domain & $R$ & Difference lower & Difference upper\\\midrule\endfirsthead
\toprule
Mapped domain & $R$ & Difference lower & Difference upper\\\midrule\endhead
Formations & 1 & 28.7755 & 28.7755\\
Formations & 1.5 & -2.2882 & 59.8019\\
Formations & 2 & -23.8622 & 81.3096\\
Formations & 3 & -52.9802 & 110.2998\\
Formations & 5 & -86.5922 & 143.7427\\
Formations & 10 & -125.3424 & 182.3383\\
Formations & 100 & -197.9409 & 255.0793\\
Formations & $\infty$ & -223.9717 & 282.0020\\
Formations + savanna & 1 & 41.8699 & 41.8699\\
Formations + savanna & 1.5 & 16.2717 & 67.4158\\
Formations + savanna & 2 & -1.5677 & 85.1567\\
Formations + savanna & 3 & -25.7930 & 109.1672\\
Formations + savanna & 5 & -54.0850 & 137.0918\\
Formations + savanna & 10 & -87.3635 & 169.7909\\
Formations + savanna & 100 & -152.9575 & 234.0673\\
Formations + savanna & $\infty$ & -178.1098 & 259.5338\\
\bottomrule
\end{longtable}


The sparse contrasts and boundary results are supplied in \path{data/density_contrast_sensitivity_F11.csv} and \path{data/density_contrast_sign_boundaries_F11.csv}; a dense $1\le R\le10$ evaluation supports the figure in \path{data/density_contrast_curve_F11.csv}. No finite contrast is adopted for the primary result. Even a justified allocation restriction would leave map/classification error, temporal correspondence and independent change validation to be addressed separately.

\section{Subsequent ecological comparisons and their admission requirements}
The broader composition-masking hypothesis requires the following ecological comparisons after measurement and reference admission:
\begin{enumerate}
\item How do independently estimated forest extent and above-ground woody biomass change on matched geographic and temporal supports?
\item When national biocapacity increases, do those physical assets increase, decrease or remain unresolved, and how much local change is offset within national totals?
\item Is account--asset divergence associated with independently observed forest-to-cropland conversion after distinguishing harvest, fire, plantation dynamics and regrowth?
\item Does adding independent asset information change an assessment when cases, weights, uncertainty and coverage are held comparable?
\end{enumerate}
These questions are not answered by conditional point-map allocation sets. Their physical-change operands require compatible forest stock and area supports; their inferential comparisons require independent paired endpoints, uncertainty including cross-year covariance and bias, reconciled territorial definitions, and appropriate inclusion and transition references. The eight national holds remain in force; account availability does not remove them. Missing field endpoints or covariance are not replaced by the density-contrast sensitivity family, spatial block arithmetic, source-access success or agreement among related map products. The finite contrast bounds identify testable auxiliary information that would be useful, but neither establish its availability nor answer the outstanding ecological questions.
